\section{Conclusion}\label{sec:conclusion}
Misinformation detectors that look near-perfect in-domain lose much of their value on a new corpus, though not all of it, and not for a single reason. Across ten audited corpora and seven detectors, \RankShareAll{} of the macro-F1 they lose is a ranking deficit that only better data or representations can address. The remaining \OpShareAll{} is a misplaced operating point, and most of it can be repaired without a single target label by transferring the share of items the detector flags rather than trusting its scores or a label-shift correction. That repair is exact when target scores are an order-preserving distortion of source scores and the base rate is unchanged, and it weakens as the base rate moves. A zero-shot language model shows the opposite balance, ranking unseen corpora better than any transferred detector while losing mostly its operating point, which the same rule repairs; published domain-generalisation methods, by contrast, do not improve on simply pooling the training corpora. Much of what in-domain benchmarks reward is the collection process behind each class: outlet signatures, attribution style, time windows. Where style is the construct itself, as with clickbait, detection travels across publishers and in part across languages; where it is a by-product of how a corpus was assembled, as with veracity, it does not. Papers that claim generalisation should report transfer, ranking metrics and the decomposition next to in-domain accuracy, and the protocol and code released with this paper make that cheap to do.
